\chapter{Literature Survey}
\label{ch:literature}

This chapter surveys the body of research directly relevant to the proposed system across five thematic areas: machine learning applied to healthcare fraud detection, explainability and interpretability in healthcare AI, evaluation of classifiers under class imbalance, clinical data validation in population health surveys, and analytical dashboards for decision support. The survey identifies the principal limitations of existing approaches and articulates the research gap that the proposed dual-domain framework addresses.

\section{Machine Learning for Healthcare Fraud Detection}
\label{sec:lit-fraud}

The application of machine learning to healthcare fraud detection has been the subject of sustained academic inquiry over the past two decades. A comprehensive systematic review by du Preez et al.\ \cite{b1} surveying supervised, unsupervised, and hybrid approaches identifies recurring challenges including limited availability of labelled training examples, the inherently evolving nature of fraud tactics, and the dependence of model performance on data quality and feature engineering choices. The review concludes that no single method dominates across all evaluation settings, and that the construction of provider-level labels from claim-level data introduces a fundamental granularity mismatch that all supervised approaches must address.

Shekhar, Leder-Luis, and Akoglu \cite{b2} investigate unsupervised machine learning for explainable healthcare fraud detection using Medicare provider data, arguing that unsupervised methods avoid the label-construction problem but sacrifice the discriminative power of supervised classifiers. Their work demonstrates that clustering-based anomaly scores are sensitive to the choice of feature normalisation and distance metric, and that the lack of ground-truth labels makes quantitative evaluation inherently difficult. The present work adopts a supervised framework grounded in the dataset-provided \textit{PotentialFraud} label while acknowledging that this label represents a provider-level aggregation rather than confirmed fraud for every associated claim.

Breiman's seminal work on Random Forests \cite{b11} establishes the theoretical basis for ensemble methods in high-dimensional classification tasks. The Random Forest algorithm constructs an ensemble of decorrelated decision trees through bootstrap aggregation and random feature subsampling, producing a robust classifier with low variance and strong performance under class imbalance when combined with balanced class weighting. Chen and Guestrin's XGBoost \cite{b12} extends gradient boosting with second-order loss approximation, regularisation, and the \texttt{scale\_pos\_weight} parameter that directly addresses the class imbalance problem in the objective function. Hosmer, Lemeshow, and Sturdivant \cite{b13} provide the canonical treatment of Logistic Regression as a linear probabilistic baseline, emphasising its interpretability and the diagnostic value of coefficient magnitudes as a first-order feature attribution mechanism. The present work systematically compares all three classifier families under identical evaluation conditions.

Provost and Fawcett \cite{b18} provide a practitioner-oriented treatment of supervised learning for business analytics that is directly applicable to healthcare claims investigation. Their emphasis on threshold selection as an operational decision --- balancing false-positive investigation costs against false-negative fraud losses --- informs the dual-threshold configuration implemented in this project. Hastie, Tibshirani, and Friedman \cite{b20} and James et al.\ \cite{b21} provide the foundational statistical learning theory underlying the cross-validation methodology and model selection framework adopted in the present work.

\section{Explainability and Interpretability in Healthcare AI}
\label{sec:lit-xai}

The deployment of machine learning in healthcare and fraud investigation contexts imposes an interpretability requirement that purely predictive performance metrics cannot satisfy. Lundberg and Lee \cite{b6} introduce SHapley Additive exPlanations (SHAP), a unified framework for explaining individual predictions of any machine learning model by computing feature contributions grounded in Shapley values from cooperative game theory. SHAP satisfies the axioms of local accuracy, missingness, and consistency --- properties that make it theoretically superior to earlier gradient-based or perturbation-based attribution methods. The \textit{TreeExplainer} variant exploits the tree structure of ensemble models to compute exact Shapley values in polynomial time, enabling on-demand per-provider explanations without batch reprocessing.

Jung et al.\ \cite{b7} conduct a systematic review of XAI properties in healthcare applications, identifying explanation fidelity, completeness, and actionability as the three properties most critical for clinical and investigative adoption. Their review finds that post-hoc global explanations alone are insufficient; investigators require local, case-specific rationales tied to individual predictions. The present work addresses this finding by implementing both global feature-importance bar charts and per-provider SHAP waterfall plots within the investigative dashboard.

Ali et al.\ \cite{b8} survey the role of explainable AI in medical and healthcare domains across 166 studies, concluding that the transparency-performance tradeoff is less severe than commonly assumed when ensemble methods with SHAP explanations are used. Their survey supports the design choice of adopting Random Forest with SHAP rather than a deep neural network with approximate attribution methods. Arrieta et al.\ \cite{b9} provide a comprehensive taxonomy of XAI techniques, categorising methods by their intrinsic or post-hoc character and their local or global scope. The taxonomy confirms that SHAP TreeExplainer occupies the post-hoc, local and global quadrant --- the most operationally relevant combination for fraud investigation workflows. Molnar \cite{b19} provides an accessible treatment of interpretable machine learning that motivates the use of both SHAP and logistic regression coefficients as complementary explanatory instruments.

\section{Evaluation of Classifiers Under Class Imbalance}
\label{sec:lit-imbalance}

Class imbalance is a structural characteristic of healthcare fraud datasets that fundamentally shapes the choice of evaluation metrics and model configuration. He and Garcia \cite{b10} provide a comprehensive review of learning from imbalanced data, cataloguing oversampling, undersampling, cost-sensitive learning, and ensemble approaches. Their analysis demonstrates that accuracy is a misleading performance metric under severe imbalance, since a classifier that predicts the majority class for all instances achieves high accuracy while providing no investigative utility. Precision, recall, F1-score, and the area under the precision-recall curve are established as the canonical metrics for imbalanced binary classification.

Chawla et al.\ \cite{b3} introduce the Synthetic Minority Over-Sampling Technique (SMOTE), which generates synthetic minority-class instances through interpolation in feature space. While SMOTE is widely adopted in the fraud detection literature, the present work opts for class-weight balancing rather than oversampling, following He and Garcia's recommendation that class-weight approaches are preferable when the minority class is small relative to feature dimensionality, as is the case with 506 positive providers across ten features.

Saito and Rehmsmeier \cite{b4} demonstrate that the precision-recall (PR) curve is more informative than the ROC curve for evaluating classifiers on imbalanced datasets, because ROC-AUC can be misleadingly optimistic when the negative class dominates. Davis and Goadrich \cite{b5} establish the theoretical relationship between PR and ROC curves, proving that a classifier that dominates in ROC space also dominates in PR space and vice versa, but that the converse does not hold under imbalance. This motivates the reporting of both ROC-AUC and Average Precision in the present evaluation --- a practice consistent with the results reported in the final research paper where Average Precision values range from 0.6611 to 0.6900 across the three classifiers.

\section{Clinical Data Validation and Population Health Surveys}
\label{sec:lit-nfhs}

The National Family Health Survey-5 \cite{b14} is the fifth in a series of nationally representative household surveys providing data on population, health, and nutrition in India. Conducted by the International Institute for Population Sciences with technical support from ICF, NFHS-5 covers 636,699 households and 724,115 female respondents across all states and union territories, collecting data on fertility, infant and child mortality, maternal and child health, and the prevalence of non-communicable diseases including diabetes and hypertension. The Ministry of Health and Family Welfare \cite{b15} publishes key indicators derived from NFHS-5 that serve as official benchmarks for national and state-level health policy.

Logical inconsistencies in self-reported survey data are a well-documented phenomenon arising from respondent recall bias, social desirability effects, interviewer error, and data entry mistakes. Clinical validation rules --- encoding medical domain knowledge as constraints over pairs of survey variables --- provide an objective, scalable mechanism for flagging implausible response combinations without requiring ground-truth annotations. The present work operationalises this approach through six weighted rules covering diabetes diagnosis-treatment inconsistencies, blood glucose threshold violations, hypertension diagnosis-medication inconsistencies, and comorbidity implausibility patterns. The weighted scoring approach, producing a continuous risk index and five-tier stratification, is novel in the NFHS-5 validation literature and extends prior binary-flag approaches by enabling priority-ordered review.

The U.S. Department of Health and Human Services Office of Inspector General \cite{b16} publishes annual reports on the Health Care Fraud and Abuse Control programme, providing the regulatory context and scale estimates that motivate the Medicare component of the present work. The CMS Medicare provider utilisation and payment data \cite{b17} constitutes the source from which the SynPUF synthetic dataset is derived, and the CMS documentation provides the authoritative definitions of claim types, reimbursement codes, and provider identification conventions used in feature engineering.

\section{Analytical Dashboards and Decision-Support Systems}
\label{sec:lit-dashboard}

Interactive analytical dashboards have become the standard mode of operationalising machine learning outputs in investigative and clinical contexts, bridging the gap between offline batch model outputs and the real-time, query-driven workflows that practitioners require. Provost and Fawcett \cite{b18} argue that the value of a predictive model is realised only through effective operationalisation, and that the design of the investigative interface --- including the presentation of uncertainty, the structure of workflow support, and the integration of model outputs with raw data --- is as important as model accuracy.

Streamlit has emerged as the dominant framework for rapid prototyping of data-centric web applications in Python, providing a declarative programming model that integrates data processing, visualisation, and interactive controls within a single script. Its multi-page application architecture, adopted in the present work, allows nine functionally distinct investigative modules to share a common state and data backend while remaining independently navigable. DuckDB \cite{b17} has attracted substantial attention as an embedded in-process analytical database that achieves columnar query performance comparable to full-server analytical systems such as Apache Spark, while requiring no separate server process and supporting direct querying of Parquet files and pandas DataFrames. Molnar \cite{b19} and Hastie et al.\ \cite{b20} collectively provide the theoretical grounding for the model evaluation and explanation visualisation components of the dashboard, including the interpretation of cross-validation learning curves and SHAP summary plots.

\section{Research Gap and Motivation for the Proposed Work}
\label{sec:lit-gap}

The literature reviewed above reveals a clear and practically significant research gap. Existing healthcare fraud detection systems operate at either the claim level or the provider level, but not both in operational linkage; no published system exposes claim-level screening evidence and provider-level ML risk scores in navigational proximity within a single investigative session. No published system combines US Medicare provider-level ML fraud detection with clinical consistency validation of Indian population health survey data within a unified platform. Existing explainable fraud detection systems apply SHAP globally or in batch mode; none integrates on-demand, per-provider SHAP waterfall plots into a live investigative dashboard. Finally, existing systems lack persistent, auditable case management that tracks investigation lifecycle status and accumulates evidence for potential legal proceedings.

The proposed system addresses each of these gaps through a principled combination of DuckDB-backed high-performance querying, multi-model supervised classification with SHAP explainability, deterministic rule-based screening, weighted clinical inconsistency validation, and SQLite-backed case management --- all integrated within a nine-page Streamlit decision-support dashboard. The design choices are grounded in and validated against the literature surveyed in this chapter.
